\documentclass[11pt]{article}
% LaTeX-only set-up (the web build does not read this file).
\usepackage[T1]{fontenc}
\usepackage[utf8]{inputenc}
\usepackage{lmodern}
\usepackage{amsmath,amssymb,amsthm}
\usepackage[a4paper,margin=1.1in]{geometry}
\usepackage{microtype}
\usepackage{booktabs}
\usepackage{enumitem}
\usepackage{newunicodechar}
\usepackage[hidelinks]{hyperref}
\setlist{itemsep=0.2em,topsep=0.4em}
% Statements carry their own explicit label through \head, so nothing is numbered automatically.
\newenvironment{statement}{\par\addvspace{\medskipamount}\noindent\ignorespaces}{\par\addvspace{\medskipamount}}
\newenvironment{frontnote}{\par\addvspace{\medskipamount}\small\noindent\ignorespaces}{\par\addvspace{\medskipamount}}
\renewcommand{\qedsymbol}{$\blacksquare$}
% Unicode that may appear inside verbatim and \texttt (Lean code)
\newunicodechar{α}{\ensuremath{\alpha}}\newunicodechar{β}{\ensuremath{\beta}}
\newunicodechar{γ}{\ensuremath{\gamma}}\newunicodechar{δ}{\ensuremath{\delta}}
\newunicodechar{μ}{\ensuremath{\mu}}\newunicodechar{ω}{\ensuremath{\omega}}
\newunicodechar{τ}{\ensuremath{\tau}}\newunicodechar{λ}{\ensuremath{\lambda}}
\newunicodechar{Π}{\ensuremath{\Pi}}\newunicodechar{Σ}{\ensuremath{\Sigma}}
\newunicodechar{∀}{\ensuremath{\forall}}\newunicodechar{∃}{\ensuremath{\exists}}
\newunicodechar{∧}{\ensuremath{\wedge}}\newunicodechar{∨}{\ensuremath{\vee}}
\newunicodechar{¬}{\ensuremath{\neg}}\newunicodechar{→}{\ensuremath{\to}}
\newunicodechar{↔}{\ensuremath{\leftrightarrow}}\newunicodechar{≤}{\ensuremath{\le}}
\newunicodechar{∈}{\ensuremath{\in}}\newunicodechar{⊢}{\ensuremath{\vdash}}
\newunicodechar{₂}{\ensuremath{_2}}\newunicodechar{₃}{\ensuremath{_3}}\newunicodechar{⁺}{\ensuremath{^+}}
% Lean listings are wider than the text font allows, so set them smaller.
\makeatletter\def\verbatim@font{\normalfont\ttfamily\footnotesize}\makeatother

% Shared by the PDF and the web build. Keep to simple \newcommand definitions.
% The short name of the main notion (alpha-plus-recursively Mahlo). Change it here and nowhere else.
% It is used both as an adjective ("the least \RM{} ordinal") and as a noun at the start of
% sentences, so it has to be an abbreviation or symbol, not a lower-case word.
\newcommand{\RM}{RM\textsuperscript{+}}
% Bold run-in head of a statement or front note: \head{Lemma 4.1}
\newcommand{\head}[1]{\textbf{#1.}}
% Named classes and predicates
\newcommand{\Mh}{\mathrm{Mh}}
\newcommand{\Sat}{\mathrm{Sat}}
\newcommand{\Mem}{\mathrm{Mem}}

\title{The least ordinal $\alpha$ that is recursively inaccessible of degree $\delta_\alpha^{k+1}$ is the least
recursively $k$-Mahlo ordinal}
% The date is stated here once; the PDF and the web page both take it from this line.
\newcommand{\paperversion}{2026-11-22}
\date{\paperversion}
\begin{document}
% Local notation (kept here so that the web build, which reads only the body, sees it too).
\newcommand{\Ad}{\mathrm{Ad}}
\newcommand{\Lim}{\mathrm{Lim}\,}
\newcommand{\CA}{C_{\mathrm{Ad}}}
\newcommand{\Lv}{\mathcal{L}}
\newcommand{\Cert}{\mathrm{Cert}}
\newcommand{\Tb}{\mathcal{T}}
\newcommand{\WF}{\mathrm{WF}}
\newcommand{\In}{\mathrm{In}}
\newcommand{\up}{{\uparrow}}
\newcommand{\Adm}{\mathrm{Adm}}
\maketitle

\begin{frontnote}\head{Abstract} For an admissible ordinal $\alpha$, let $\delta_\alpha$ be Aguilera's supremum of the
order types of well-orderings $\Delta_1$-definable over $L_\alpha$ with parameters. An ordinal $\alpha$ is recursively
inaccessible of degree $\delta_\alpha$ if it lies in every class obtained by iterating ``admissible limit of'' along a
node in the well-founded part of a $\Sigma_1$ presentation at $\alpha$; degree $\delta_\alpha^{k+2}$ iterates the
passage to degree $\delta_\alpha^{k+1}$ in the same way. The recursively $0$-Mahlo ordinals are the recursively Mahlo
ones, and the recursively $(k+1)$-Mahlo ordinals are those $\Pi_2$-reflecting on the recursively $k$-Mahlo ordinals.
Up to and including the least (+1)-stable ordinal, the admissible ordinals of degree $\delta_\alpha^{k+1}$ are exactly
the recursively $k$-Mahlo ordinals and their admissible limits, so the least of them is the least recursively
$k$-Mahlo ordinal. The correspondence stops at the first limit: the least ordinal of degree $\delta_\alpha^\omega$ is
below the least recursively $\omega$-Mahlo ordinal.
\end{frontnote}

\begin{frontnote}\head{Provenance} The proofs were found and written by AI systems. No human specialist has
reviewed them, nothing is formalised in Lean, and novelty is not established. Section 9 states what has been checked.
\end{frontnote}

\section{Introduction}
% key-steps: Lemma 4.4, Theorem 6.5

We build on [A] and [E] (see the references) and use their conventions and numbering, citing their items as [A,
Theorem 5.7] or [E, (C2)]. [A] shows that the least ordinal $\alpha$ that is recursively inaccessible of degree
$\delta_\alpha$ is the least recursively Mahlo ordinal, and that up to and including the least (+1)-stable
ordinal the admissible ordinals of that degree are the recursively Mahlo ordinals and their
admissible limits [A, Corollary 5.9]. Its third open question asks what happens at higher degrees.

\textbf{Degrees above $\delta_\alpha$.} At admissible $\alpha > \omega$, the ranks of nodes of $\Sigma_1$ presentations are exactly
the ordinals below $\delta_\alpha$ [A, Theorem 3.1], so higher degrees cannot be ranks of such nodes; the names below record
how a class is built, as in ordinal arithmetic, not a rank. Degree $\delta_\alpha + 1$ is ``admissible and a limit of
the class of degree $\delta_\alpha$''. Degree $\delta_\alpha \cdot 2$ is $\delta_\alpha$ further steps of ``admissible
limit of'' over the class of degree $\delta_\alpha$. Degree $\delta_\alpha^{k+2}$ asks, for every well-founded node of
every $\Sigma_1$ presentation, for membership in the iteration along that node of the passage to degree
$\delta_\alpha^{k+1}$ (Definitions 3.1--3.4).

The main result is the following.

\begin{statement}\head{Theorem 7.2 (main theorem)} Let $k \ge 0$. An admissible ordinal $\gamma$ that is at most
the least (+1)-stable ordinal is recursively inaccessible of degree $\delta_\gamma^{k+1}$ if and only if $\gamma$ is
recursively $k$-Mahlo or a limit of recursively $k$-Mahlo ordinals. The least ordinal $\alpha$ that is recursively
inaccessible of degree $\delta_\alpha^{k+1}$ is the least recursively $k$-Mahlo ordinal.
\end{statement}

Recursive $k$-Mahloness is defined in Section 2.3. For
$k = 0$ this is [A, Corollary 5.9]; for $k = 1$, degree $\delta_\alpha^2$ gives the least recursively 1-Mahlo ordinal.
The theorem is the case $X = \Ad$ of Theorem 7.1, which describes the classes $X[k]$ of degree $\delta^k$ over any base
class $X$ built in Section 3: at admissible $\gamma$ at most the least (+1)-stable ordinal, $X[k]$ consists of the members and admissible
limits of the ordinals in $\Lim X$ that are admissible if $k = 0$ and recursively $(k-1)$-Mahlo if $k \ge 1$.

\textbf{Further results} (Section 7).
\begin{itemize}
\item \emph{Degrees $\delta_\alpha + 1$ and $\delta_\alpha \cdot 2$} (Corollary 7.3). The least ordinal of degree
  $\delta_\alpha + 1$ is the least admissible limit of recursively Mahlo ordinals. The least ordinal of degree
  $\delta_\alpha \cdot 2$ is the least recursively Mahlo limit of recursively Mahlo ordinals, which is below the least
  recursively 1-Mahlo ordinal.
\item \emph{Degree $\delta_\alpha^\omega$} (Proposition 7.4). For admissible $\gamma$ at most the least (+1)-stable ordinal, degree
  $\delta_\gamma^\omega$ holds exactly when $\gamma$ is a limit of recursively $k$-Mahlo ordinals for every $k <
  \omega$. Its least element, the least admissible ordinal with this property, is below the least recursively
  $\omega$-Mahlo ordinal.
\end{itemize}
Section 8 says where the statements fail: at exponent $0$, at degree $\delta_\alpha \cdot 2$, at exponent $\omega$, and
for step classes tested with a parameter.

\textbf{How the proof goes.} The upper bound (Theorem 5.2) is the argument of [A, Theorem 5.7(a)] one level up: a
member of $\Lim X$ that is admissible if $k = 0$ and recursively $(k-1)$-Mahlo if $k \ge 1$ reflects the $\Pi_2$ sentences saying that the classes required of it are
unbounded, and so lies in $X[k]$. The lower bound (Theorem 6.5) is by induction on $k$ for all base classes at once. If
no recursively $k$-Mahlo member of $\Lim X$ lies in $(\ell, \gamma]$, a $\Sigma_1$ tree searching for such a point is
well-founded at $\gamma$. The class $X[k+1]$ asks for membership in the iteration of $Y \mapsto Y[k]$ along this tree;
applying the induction hypothesis to the base classes of that iteration, a least-counterexample descent excludes
$\gamma$. This needs membership in every class $X[k]$ to be decided by a bounded test on the constructible levels up to
the next admissible ordinal (Lemma 4.4), because those base classes must be tested inside the $\Sigma_1$ tree.
Appendix A lists which result uses which.

\textbf{Nearest results.} Richter and Aczel (Theorems 1.9--1.11) show that recursively Mahlo ordinals lie in the
iterates of ``admissible limit of'' considered there; [A, Theorem 5.7(a)] gives the upper bound for nodes in the
externally well-founded part of $\Sigma_1$ presentations, which is the upper inclusion of Theorem 7.2 at $k = 0$, obtained here from Theorem 5.2(A) with $X = \mathrm{Ad}$ and $c = 1$. Iterations of the Mahlo operation
along orderings, with ordinals or notations as indices, appear in Arai, Lectures on Ordinal Analysis (Section 4, one
level higher, inside $\Pi_3$-reflection) and Iterating the recursively Mahlo operations (Sections 2--3), and in Rathjen,
Proof theory of reflection (Remark 4.9, for cardinals); we cite them for comparison only. These works give iterations
and upper bounds. We did not find in them degrees beyond $\delta_\alpha$ defined through $\Sigma_1$ presentations, a
lower bound of the kind proved here, or the failure at $\omega$.

\section{Preliminaries}

\subsection{\texorpdfstring{Conventions from [E] and [A]}{Conventions from [E] and [A]}}

We use [A, Section 2] without change; it collects the conventions of [E]. In brief:
\begin{itemize}
\item \emph{Reflection.} An ordinal $\mu$ is \emph{$\Pi_n$-reflecting on} a class $Y$ if $\mu$ is a limit above
  $\omega$ and every $\Pi_n$ sentence $\varphi(q)$ with $q \in L_\mu$ true in $L_\mu$ is true in $L_\beta$ for some
  $\beta \in Y \cap \mu$ with $q \in L_\beta$; with $Y$ the class of all ordinals, $\mu$ is \emph{$\Pi_n$-reflecting}.
  A limit $\mu > \omega$ is $\Pi_2$-reflecting if and only if it is admissible (Richter--Aczel, Theorem 1.8). In a limit
  level a block of like quantifiers can be contracted to one quantifier over tuples (\emph{block contraction}), so there
  every $\Pi_2$ sentence is equivalent to one of the form $\forall u\, \exists v\, \theta(q,u,v)$ with $\theta$ a
  $\Delta_0$ formula, and a conjunction of $\Pi_2$ sentences is a $\Pi_2$ sentence. $\mathrm{Ord}(y)$ says that $y$ is
  an ordinal.
\item \emph{Classes.} $\Lim Y$ is the class of limit ordinals $\mu$ with $Y \cap \mu$ unbounded in $\mu$. $\Ad$ is the
  class of admissible ordinals above $\omega$, and $\mu^+$ is the least admissible ordinal above $\mu$. Recursively
  inaccessible ordinals are those in $\Ad \cap \Lim \Ad$.
\item \emph{Certificates} [E, Section 2]. $\Cert(W, \xi)$ is a $\Delta_0$ formula saying that $W$ is a
  \emph{certificate through} $\xi$: a set consisting of [E]'s syntax data and, for each $\nu \le \xi$, a set
  $\Lv^W_\nu$ and a truth table $\mathrm{Tr}^W_\nu$ (we drop the superscript when the certificate is clear). (C1): every
  certificate is the true one, $\Lv_\nu = L_\nu$ and $\mathrm{Tr}_\nu$ the satisfaction relation of $L_\nu$; a pair
  $(\varphi, \langle q\rangle)$ is in $\mathrm{Tr}_\nu$ when $L_\nu \models \varphi(q)$, and a sentence is entered with
  the dummy constant $\emptyset$. (C2): an admissible $L_\alpha$, $\alpha > \omega$, contains a certificate through
  every $\xi < \alpha$. (C3): the $\Pi_2$ sentence $G$ (``$\omega$ exists, every ordinal has a successor, every ordinal
  has a certificate'') holds at admissible levels above $\omega$, and a nonempty level satisfying it is a limit level
  above $\omega$. (C4): $\Sat_2(\varphi, q)$ is a $\Pi_2$ formula that, at levels satisfying $G$, holds exactly when
  $\varphi$ is the code of a $\Pi_2$ formula $\forall u\, \exists v\, \theta$ and $\varphi(q)$ is true. (C5): $\Sigma_1$
  sentences persist upward. For $\nu \le \xi$, $W{\restriction}(\nu+1)$ is $W$ cut down to its rows of index at most
  $\nu$; it is used only inside bounded conditions, where quantifiers over it are read as quantifiers over $W$ guarded by
  the index [A, Section 2.3].
\item \emph{Codes.} Formula codes are hereditarily finite sets. A primitive recursive relation on codes (being the
  code of a $\Sigma_1$ or $\Pi_2$ formula, the enumeration of tasks in Section 6) is defined over $L_\omega$ by a fixed
  formula, so it is read from $\mathrm{Tr}_\omega$ of any certificate through at least $\omega$ by a bounded condition [A, Section 2.1].
\item \emph{Countability} [E, Lemma 7.2]. For every admissible $\gamma$ at most the least (+1)-stable ordinal
  (Table 1), every $L_\eta$ with $\eta < \gamma$ is countable in $L_\gamma$ [A, Section 2.2].
\item \emph{Background.} $\Delta_0$ absoluteness, and $\Sigma$-recursion, $\Delta_0$-separation and
  $\Delta_0$-collection in admissible levels (Barwise, Corollary I.8.6 and Theorem I.6.4), as in [E, Section 2].
\end{itemize}

\subsection{\texorpdfstring{Presentations and the degree $\delta_\alpha$}{Presentations and the degree delta alpha}}

Presentations are as in [A, Section 2.3]. A \emph{presentation} is a pair $Q = (\rho_Q, p_Q)$ of a formula
$\rho_Q(x,y,z)$ and a set $p_Q$, its \emph{parameter}. At a limit $\mu > \omega$ with $p_Q \in L_\mu$, the elements of
$L_\mu$ are \emph{nodes}, and $e'$ is a \emph{predecessor} of $e$ at $\mu$ if $L_\mu \models \rho_Q(e', e, p_Q)$; the
pair $(e', e)$ is then an \emph{edge} at $\mu$. $Q$ is a \emph{$\Sigma_1$ presentation} or a \emph{$\Delta_0$
presentation} if $\rho_Q$ is $\Sigma_1$ or $\Delta_0$. $\WF(Q, \mu)$, the \emph{well-founded part} at $\mu$, is the set
of nodes from which there is no infinite sequence $e_0, e_1, \dots$ with $e_{i+1}$ a predecessor of $e_i$ at $\mu$,
taken in the universe, and $\mathrm{rank}_\mu(e)$ is the rank of $e$ there. By (C5) an edge of a $\Sigma_1$ presentation
at $\mu$ is an edge at every larger level, so if $e \in \WF(Q,\gamma)$ and $\mu < \gamma$ is a limit with $p_Q, e \in
L_\mu$, then $e \in \WF(Q,\mu)$. A \emph{set relation} is a $\Delta_0$ presentation whose edges are the same at every
level containing its parameter.

An \emph{eligible pair at $\mu$} is a pair $(Q, e)$ of a $\Sigma_1$ presentation $Q$ with $p_Q \in L_\mu$ and a node $e
\in \WF(Q, \mu)$.

For a limit $\alpha > \omega$, $\delta_\alpha$ is the supremum of the order types of well-orderings with field a
subclass of $L_\alpha$ that are $\Delta_1$-definable over $L_\alpha$ with parameters (Aguilera, The order of
reflection, Section 2); for admissible $\alpha$ these are exactly the ranks of the nodes in $\WF(Q,\alpha)$ for
eligible pairs $(Q, e)$ at $\alpha$ [A, Theorem 3.1]. By [A, Definition 2.1], $\mu \in \CA(Q, e)$ if and only if $\mu
\in \Ad$, $p_Q, e \in L_\mu$, and $\mu \in \Lim \CA(Q, e')$ for every predecessor $e'$ of $e$ at $\mu$. A limit $\alpha
> \omega$ is \emph{recursively inaccessible of degree $\delta_\alpha$} if $\alpha \in \CA(Q, e)$ for every eligible pair
$(Q,e)$ at $\alpha$ [A, Definition 2.2]. We write $D$ for the class of these ordinals.

\subsection{The finite Mahlo classes}

The \emph{recursively Mahlo} ordinals are those $\Pi_2$-reflecting on $\Ad$. As in [A], the ordinals
$\Pi_2$-reflecting on the recursively Mahlo ordinals are called \emph{recursively 1-Mahlo}, just as ``1-inaccessible''
is one step above ``inaccessible''; in general the ordinals $\Pi_2$-reflecting on the recursively $k$-Mahlo ordinals
are \emph{recursively $(k+1)$-Mahlo}, and recursively 0-Mahlo means recursively Mahlo. The ordinals that are
recursively $k$-Mahlo for every $k$ are \emph{recursively $\omega$-Mahlo}. [E, Section 3.1] defines the same hierarchy
starting from the limits above $\omega$, so that its index is $k+2$.

In formulas and proofs we use the $\Pi_2$-iterates. For a class $X$, $\Pi_2(X)$ is the class of ordinals
$\Pi_2$-reflecting on $X$, and $\Pi_2$ alone is the class of $\Pi_2$-reflecting ordinals, that is, $\Ad$. Let
$\Pi_2^0$ be the class of limit ordinals above $\omega$; being $\Pi_2$-reflecting on it is the same as being
$\Pi_2$-reflecting, so $\Pi_2^1 = \Pi_2(\Pi_2^0) = \Pi_2$. Put $\Pi_2^{\eta+1} = \Pi_2(\Pi_2^\eta)$, and $\Pi_2^\lambda = \bigcap_{\eta<\lambda} \Pi_2^\eta$ at
limits. So $\Pi_2^2$ is the class of recursively Mahlo ordinals, $\Pi_2^{k+2}$ that of the recursively $k$-Mahlo
ordinals, and $\Pi_2^\omega$ that of the recursively $\omega$-Mahlo ordinals. Theorem 7.1 matches degree $\delta^k$
with $\Pi_2^{k+1}$.

\begin{center}\begin{tabular}{p{0.14\linewidth}p{0.76\linewidth}}
\toprule
Ordinal & Definition \\
\midrule
$m_\omega$ & the least recursively $\omega$-Mahlo ordinal \\
$\pi_3$ & the least $\Pi_3$-reflecting ordinal \\
(+1)-stable & $\alpha$ is \emph{(+1)-stable} if $L_\alpha \prec_{\Sigma_1} L_{\alpha+1}$; equivalently, $\alpha$ is
$\Pi_n$-reflecting for every finite $n$ (Richter--Aczel, Theorem 1.18). The least (+1)-stable ordinal exists by [E,
proof of Lemma 7.1] \\
\bottomrule
\end{tabular}\\[0.4em]Table 1. Ordinals used below.\end{center}

\begin{statement}\head{Lemma 2.1} (a) For every $k < \omega$, every recursively $(k+1)$-Mahlo ordinal is
recursively $k$-Mahlo and a limit of recursively $k$-Mahlo ordinals, every recursively Mahlo ordinal is admissible
and a limit of admissible ordinals, and every admissible ordinal above $\omega$ is a limit above $\omega$ and a limit
of limits above $\omega$; in particular every recursively $k$-Mahlo ordinal is a limit of admissible
ordinals.
(b) $\pi_3$ is $\Pi_2$-reflecting on the recursively $\omega$-Mahlo ordinals. Hence $m_\omega$ and the least
admissible ordinal that is a limit of recursively $k$-Mahlo ordinals for every $k < \omega$ exist, the latter is at
most $m_\omega$, $m_\omega < \pi_3$, and $\pi_3$ is at most the least (+1)-stable ordinal.
\end{statement}

\begin{proof}
(a) By induction on $k$, since $\Pi_2$-reflection on a class implies it on every larger class; the second inclusion
follows by reflecting the sentence ``$y$ exists'' for each $y < \mu$.

(b) By [A, Theorem 4.3] with $n = 2$, $\pi_3$ lies in [E]'s class $C(Q, e)$ [E, Definition 3.1] for every eligible
pair $(Q,e)$ at $\pi_3$. For the set relation $\in$ on $\omega+2$ and its node $\omega+1$, of rank $\omega+1$, this class
is the class of ordinals $\Pi_2$-reflecting on the recursively $\omega$-Mahlo ordinals [E, Example 3.2(1)], and
reflecting a trivial sentence gives a recursively $\omega$-Mahlo ordinal below $\pi_3$. The ordinal $m_\omega$ is
admissible and, by (a), a limit of recursively $k$-Mahlo ordinals for every $k$, so the least admissible such ordinal is
at most $m_\omega$. Finally the least (+1)-stable ordinal is $\Pi_3$-reflecting, so $\pi_3$ is at most it.
\end{proof}

\section{\texorpdfstring{Degrees above $\delta_\alpha$}{Degrees above delta alpha}}

\subsection{One further step}

\begin{statement}\head{Definition 3.1} Let $X$ be a class of ordinals and $Q$ a presentation. By recursion on $\mu$,
$\mu \in \In_X(Q, e)$ if and only if $\mu \in \Ad \cap \Lim X$, $p_Q, e \in L_\mu$, and $\mu \in \Lim \In_X(Q, e')$ for
every predecessor $e'$ of $e$ at $\mu$. Let $\up X$ be the class of limits $\mu > \omega$ with $\mu \in \In_X(Q, e)$ for
every eligible pair $(Q, e)$ at $\mu$.

An ordinal $\alpha$ is \emph{recursively inaccessible of degree $\delta_\alpha + 1$} if $\alpha \in \Ad \cap \Lim D$, and
\emph{of degree $\delta_\alpha \cdot 2$} if $\alpha \in \up D$.
\end{statement}

The class $\In_X(Q,e)$ is $\CA(Q,e)$ with ``$\mu \in \Lim X$'' added, so $\up X$ is $\delta_\alpha$ further steps of
``admissible limit of'' over $X$. Every class $\In_X(Q,e)$ is contained in $\Ad \cap \Lim X$, and so is $\up X$: the
$\Delta_0$ presentation with parameter $\emptyset$ and no edges gives an eligible pair whose class is $\Ad \cap \Lim X$.

\subsection{Class terms}

The exponent of a degree is a node of the chain $\omega + 1$: let $P_\omega$ be the $\Delta_0$ presentation ``$x \in y
\wedge y \in z$'' with parameter $\omega+1$, whose node $b \le \omega$ has the ordinals $c < b$ as predecessors at every
level containing $\omega + 1$.

\emph{Example.} Over a base class $X$, the class $X[1]$ is $\up X$, and $X[2]$ is defined as follows. For an eligible
pair $(Q, e)$ let $H_{X,1}(Q, e)$ be the class of the $\mu \in \Ad \cap \Lim X$ (with $p_Q, e \in L_\mu$ and $e \in
\WF(Q,\mu)$) that lie in $\up H_{X,1}(Q, e')$ for every predecessor $e'$ of $e$ at $\mu$; so $H_{X,1}(Q, \cdot)$
iterates $Y \mapsto \up Y$ along $Q$, starting from $\Ad \cap \Lim X$ at the nodes without predecessors. Then $X[2]$
consists of the $\mu \in \up X$ that lie in $H_{X,1}(Q,e)$ for every eligible pair $(Q,e)$ at $\mu$. In general $X[b]$
asks, for each $c < b$ and each eligible pair $(Q,e)$, for membership in the iteration $H_{X,c}(Q, \cdot)$ of $Y \mapsto
Y[c]$ along $Q$. Since $\up$ itself quantifies over all eligible pairs, the classes are named by finite expressions and
defined by one simultaneous recursion.

\begin{statement}\head{Definition 3.2 (class terms)} \emph{Terms} are built by three rules: $\Ad$ is a term; if $X$ is a
term and $b \le \omega$, then $X[b]$ is a term; if $X$ is a term, $c < \omega$, $Q = (\rho_Q, p_Q)$ is a $\Sigma_1$
presentation and $e$ is a set, then $H_{X,c}(Q, e)$ is a term. In the last two, $X$ is the \emph{base term}. Membership
of an ordinal $\mu$ in a term is defined by recursion on $\mu$:
\begin{itemize}
\item $\mu \in \Ad$ if $\mu$ is admissible and above $\omega$;
\item $\mu \in X[b]$ if and only if $\mu \in \Ad \cap \Lim X$ and $\mu \in H_{X,c}(Q, e)$ for every $c < b$ and every
  eligible pair $(Q,e)$ at $\mu$;
\item $\mu \in H_{X,c}(Q, e)$ if and only if $\mu \in \Ad \cap \Lim X$, $p_Q, e \in L_\mu$, $e \in \WF(Q, \mu)$, and
  $\mu \in H_{X,c}(Q, e')[c]$ for every predecessor $e'$ of $e$ at $\mu$.
\end{itemize}
The \emph{constants} of a term are the parameters $p_Q$ and nodes $e$ in it, together with the constants of its base
term. Its \emph{code} $t_X$ is $\langle 0\rangle$ for $\Ad$, $\langle 1, t_X, b\rangle$ for $X[b]$, and $\langle 2, t_X,
c, \rho_Q, p_Q, e\rangle$ for $H_{X,c}(Q,e)$.
\end{statement}

A set is a \emph{term code} if some finite sequence ending with it starts with $\langle 0\rangle$ and forms each later
entry from an earlier one by the second or third rule, the formula codes being codes of $\Sigma_1$ formulas. The
entries are codes of subterms, so for a limit $\lambda$ and a code in $L_\lambda$ such a sequence lies in $L_\lambda$;
being a term code is therefore a bounded condition on a certificate through $\lambda$, with the search bounded by
$\Lv_\lambda$.

\begin{statement}\head{Lemma 3.3} (a) The clauses of Definition 3.2 determine the membership of every ordinal in every
term. (b) Every term is a subclass of $\Ad$, and if $\mu \in X$ then every constant of $X$, and the code $t_X$, lies
in $L_\mu$.
\end{statement}

\begin{proof}
(a) By recursion on $\mu$, for all terms at once; assume membership at every $\mu' < \mu$ is determined, so that $\Lim Y$
at $\mu$ is determined for every term $Y$. A statement $\mu \in H_{Y,c}(Q,e)$ whose guards fail is false. Otherwise
give ``$\mu \in Y[b]$'' the measure $(b, 0, 0)$ and ``$\mu \in H_{Y,c}(Q,e)$'' the measure $(c, 1,
\mathrm{rank}_\mu(e))$, ordered lexicographically. The clause of $Y[b]$ refers, besides $\Lim Y$, to statements with
first coordinate $c < b$, and the clause of $H_{Y,c}(Q,e)$ refers, besides $\Lim Y$, to statements $\mu \in Y'[c]$ with
measure $(c, 0, 0)$. So recursion on the measure determines membership at $\mu$.

(b) Every clause requires $\mu \in \Ad$. For the constants, induction on the term: the clause of $H_{X,c}(Q,e)$ requires
$p_Q, e \in L_\mu$, and $\mu \in \Lim X$ gives some $\eta < \mu$ in $X$, whose constants are in $L_\eta \subseteq L_\mu$.
Formula codes and $b, c \le \omega$ are in $L_\mu$, and $L_\mu$ is closed under finite tuples.
\end{proof}

\begin{statement}\head{Definition 3.4} For $0 \le k < \omega$, an ordinal $\alpha$ is \emph{recursively inaccessible of
degree $\delta_\alpha^{k+1}$} if $\alpha \in \Ad[k+1]$, and \emph{of degree $\delta_\alpha^\omega$} if $\alpha \in
\Ad[\omega]$.
\end{statement}

\begin{statement}\head{Lemma 3.5} Let $X$ be a term and $Q$ a $\Sigma_1$ presentation.
\begin{itemize}
\item[(a)] If $p_Q, e \in L_\mu$ and $e \in \WF(Q, \mu)$, then $\mu \in H_{X,0}(Q,e)$ if and only if $\mu \in
  \In_X(Q,e)$.
\item[(b)] $X[1] = \up X$.
\item[(c)] $\Ad[1] = D$: an ordinal is recursively inaccessible of degree $\delta_\alpha^1$ exactly when it is
  recursively inaccessible of degree $\delta_\alpha$.
\end{itemize}
\end{statement}

\begin{proof}
(a) For every term $Y$, $Y[0] = \Ad \cap \Lim Y$, since $0$ has no predecessors. Hence $\mu \in H_{X,0}(Q,e)$ if and only
if $\mu \in \Ad \cap \Lim X$, $p_Q, e \in L_\mu$, and $\mu \in \Lim H_{X,0}(Q,e')$ for every predecessor $e'$ of $e$ at
$\mu$. This is the clause of $\In_X(Q,e)$ with $H_{X,0}(Q,e')$ in place of $\In_X(Q,e')$. By induction on $\mu$ the two
classes agree: at a level $\mu' < \mu$ with $p_Q, e' \in L_{\mu'}$, $e' \in \WF(Q, \mu')$ by (C5), so the induction
hypothesis applies at $\mu'$, and levels not containing $e'$ are in neither class.

(b) The only predecessor of $1$ is $0$, so by (a) $\mu \in X[1]$ if and only if $\mu \in \Ad \cap \Lim X$ and $\mu \in
\In_X(Q,e)$ for every eligible pair $(Q,e)$ at $\mu$. Since every class $\In_X(Q,e)$ is contained in $\Ad \cap \Lim X$,
this is $\mu \in \up X$.

(c) By (b), $\Ad[1] = \up \Ad$. By induction on $\mu$, $\In_\Ad(Q,e) \subseteq \CA(Q,e)$, so $\up \Ad \subseteq D$. For
the converse let $Q^+$ be the presentation with the same parameter whose nodes are the pairs $\hat x = (0, x)$ and
$\bot = (1, 0)$, with edges $(\hat x, \hat y)$ for the edges $(x, y)$ of $Q$ and $(\bot, \hat y)$ for every $y$. It is
$\Sigma_1$, and $\hat e \in \WF(Q^+, \mu)$ whenever $e \in \WF(Q,\mu)$. The class $\CA(Q^+, \bot)$ is $\Ad$ restricted to
the levels containing $p_Q$, so by induction on $\mu$, $\CA(Q^+, \hat e) = \In_\Ad(Q, e)$. Hence every member of $D$ is
in $\up \Ad$.
\end{proof}

By Lemma 3.5, $D = \Ad[1]$ is a term and the class of degree $\delta_\alpha \cdot 2$ is $D[1]$. The class $\Ad[0] = \Ad
\cap \Lim \Ad$ consists of the recursively inaccessible ordinals.

\begin{center}\begin{tabular}{p{0.22\linewidth}p{0.68\linewidth}}
\toprule
Symbol & Meaning \\
\midrule
$\In_X(Q,e)$, $\up X$ & iteration of ``admissible limit of'' over $X$ along $e$; along all eligible pairs (Def.\ 3.1) \\
$D$ & the ordinals recursively inaccessible of degree $\delta_\alpha$ (Section 2.2); $D = \Ad[1]$ \\
$X[b]$ & degree $\delta^b$ over the base term $X$, for $b \le \omega$ (Def.\ 3.2) \\
$H_{X,c}(Q,e)$ & iteration of $Y \mapsto Y[c]$ along $e$, over $X$ (Def.\ 3.2) \\
$t_X$ & the code of the term $X$ \\
\bottomrule
\end{tabular}\\[0.4em]Table 2. Classes of Section 3.\end{center}

\section{Bounded tests one admissible level up}

Membership in a term at $\nu$ quantifies over the well-founded part at $\nu$, which is not decided inside $L_\nu$. This
section shows that it is decided by a bounded test on a certificate through $\nu^+$, uniformly in the code of the term
(Lemma 4.4), and gives a bounded test, uniform in $k$, for the classes $\Pi_2^{k+1}$ (Lemma 4.1).

\begin{statement}\head{Lemma 4.1} There is a $\Delta_0$ formula $\chi(W, \nu, k)$ without parameters such that for all
ordinals $\nu$, all $k < \omega$ and every certificate $W$ through $\nu$: $\chi(W,\nu,k)$ holds if and only if $\nu$
is a limit above $\omega$ (for $k = 0$), admissible (for $k = 1$), or recursively $(k-2)$-Mahlo (for $k \ge 2$).
\end{statement}

\begin{proof}
The class in the statement is $\Pi_2^k$, and $\Pi_2^{k+1}$ is the class of ordinals $\Pi_2$-reflecting on
$\Pi_2^k$. Let $F_0$ be the sentence ``$\omega$ exists and every ordinal has a successor''; for every ordinal $\eta$, $L_\eta
\models F_0$ if and only if $\eta$ is a limit above $\omega$. Let
\[ \begin{aligned} F_{k+1} = {} & F_0 \wedge G \wedge \forall \varphi\, \forall q\, [\, \neg\Sat_2(\varphi, q) \vee
   \exists \eta\, \exists W' (\Cert(W', \eta) \wedge q \in \Lv_\eta \\ & \wedge (\varphi, \langle q\rangle) \in
   \mathrm{Tr}_\eta \wedge (\ulcorner F_k\urcorner, \langle \emptyset\rangle) \in \mathrm{Tr}_\eta)\, ], \end{aligned} \]
where $\ulcorner F_k\urcorner$ is named by a parameter-free $\Delta_0$ definition of that hereditarily finite set. The
map $k \mapsto \ulcorner F_k\urcorner$ is primitive recursive, so its graph is defined over $L_\omega$ by a fixed formula
$\sigma$. Put
\[ \chi(W, \nu, k) = \omega \le \nu \wedge \exists f \in \Lv_\omega\, [(\sigma, \langle k, f\rangle) \in
   \mathrm{Tr}_\omega \wedge (f, \langle \emptyset\rangle) \in \mathrm{Tr}_\nu]. \]
It is bounded. By (C1) it says that $L_\nu \models F_k$, so it suffices to show, by induction on $k$, that $L_\nu
\models F_k$ if and only if $\nu \in \Pi_2^k$. The complexity of $F_k$ plays no role, since its truth is read from the
truth table.

For $k = 0$ this is the property of $F_0$. Suppose it holds for $k$. If $L_\nu \models F_{k+1}$, then $\nu$ is a limit
above $\omega$ and $L_\nu \models G$. Let $\varphi(q)$ be a $\Pi_2$ sentence true in $L_\nu$, in the form $\forall u\,
\exists v\, \theta$. By (C4), $L_\nu \models \Sat_2(\varphi, q)$, so there are $\eta < \nu$ and a certificate $W' \in
L_\nu$ through $\eta$, true by (C1), with $q \in L_\eta$, $L_\eta \models \varphi(q)$ and $L_\eta \models F_k$, that is,
$\eta \in \Pi_2^k$. So $\nu \in \Pi_2^{k+1}$. Conversely let $\nu \in \Pi_2^{k+1}$. Then $\nu$ is $\Pi_2$-reflecting, hence
admissible (Richter--Aczel 1.8), so $L_\nu \models F_0 \wedge G$ by (C3). If $L_\nu \models \Sat_2(\varphi, q)$, then
$\varphi(q)$ is a $\Pi_2$ sentence true in $L_\nu$ (C4), so some $\eta \in \Pi_2^k \cap \nu$ has $q \in L_\eta$ and
$L_\eta \models \varphi(q)$; also $L_\eta \models F_k$, and a certificate through $\eta$ is in $L_\nu$ by (C2).
\end{proof}

For a certificate $V$ through $\eta$ and $\kappa \le \eta$ write $\Adm(V, \kappa)$ for $\chi(V{\restriction}(\kappa+1),
\kappa, 1)$; by Lemma 4.1 it says that $\kappa \in \Ad$. ``$\kappa = \nu^+$ according to $V$'' is the bounded condition
$\nu < \kappa \wedge \Adm(V,\kappa)$ ${}\wedge \forall \kappa' < \kappa\, (\nu < \kappa' \rightarrow \neg\Adm(V,\kappa'))$,
and ``$\nu^+ \le \eta$ according to $V$'' is $\exists \kappa \le \eta\, (\nu < \kappa \wedge \Adm(V,\kappa))$. The
ordinal $\eta$ through which $V$ is a certificate is read off from $V$ (its largest row index), so it is not a separate
argument of the formulas below.

\begin{statement}\head{Lemma 4.2 (well-foundedness one level up)} Let $\lambda > \omega$ be a limit, $Q = (\rho_Q, p_Q)$
a presentation with $p_Q \in L_\lambda$, $e \in L_\lambda$, and $V$ a certificate through some $\eta \ge \lambda^+$.
Let $\WF_V(Q, \lambda, e)$ be the condition: there are $\kappa \le \eta$ with $\kappa = \lambda^+$ according to $V$,
and $Z, g \in \Lv_\kappa$ with $e \in Z \subseteq \Lv_\lambda$, $g$ a function from $Z$ to ordinals, and for all $a \in
Z$ and $e' \in \Lv_\lambda$ with $(\rho_Q, \langle e', a, p_Q\rangle) \in \mathrm{Tr}_\lambda$: $e' \in Z$ and $g(e')
< g(a)$. Then $\WF_V(Q,\lambda,e)$ is $\Delta_0$ in $V, p_Q, e$, and it holds if and only if $e \in \WF(Q, \lambda)$.
\end{statement}

\begin{proof}
By (C1) $V$ is the true certificate, so $\Lv_\kappa = L_{\lambda^+}$ and the edges read from $\mathrm{Tr}_\lambda$ are
the edges of $Q$ at $\lambda$.

\emph{Sufficiency.} An infinite descending sequence from $e$ stays in $Z$, and $g$ would map it to a strictly
decreasing sequence of ordinals.

\emph{Necessity.} The edge relation $E = \{(e', a) : L_\lambda \models \rho_Q(e', a, p_Q)\}$ is definable over
$L_\lambda$, and $L_\lambda$ is closed under pairs, so $E \in L_{\lambda+1} \subseteq L_{\lambda^+}$. Call $(Z, g)$
\emph{good for $a$} if $a \in Z \subseteq L_\lambda$, $Z$ is closed under $E$-predecessors, and $g$ is the exact rank
function on $Z$: $g(x) = \sup\{g(y) + 1 : y \mathrel{E} x\}$ for $x \in Z$, the empty supremum being $0$. Given $E$, $Z$
and $g$ this is a bounded condition. Two good pairs agree on the intersection of their domains, which is closed under
predecessors and contained in $\WF(Q,\lambda)$, because exact rank functions on such a set are unique. We show by
induction along $E$ on $\WF(Q, \lambda)$, in the universe, that every $a \in \WF(Q,\lambda)$ has a good pair in
$L_{\lambda^+}$. All predecessors of $a$ are in $\WF(Q,\lambda)$, and the set of them is in $L_{\lambda^+}$ by
$\Delta_0$-separation. By the induction hypothesis, $L_{\lambda^+}$ satisfies ``every predecessor of $a$ has a good
pair'', so $\Delta_0$-collection in the admissible $L_{\lambda^+}$ gives a set of good pairs covering the predecessors.
Their union is a function, by agreement, on a set closed under predecessors; adding $a$ with the supremum of $g(e') + 1$
over its predecessors $e'$ gives a good pair for $a$. An exact rank function decreases along edges, so a good pair for
$e$ satisfies the condition.
\end{proof}

\textbf{The test at one level.} Fix an admissible $\lambda$, the function $\Tb$ on $\lambda$ with $\Tb(\eta) = \{t \in
L_\eta : t \text{ is a term code and } \eta \in t\}$ (the \emph{table of terms below $\lambda$}), and a certificate $V$
through some $\eta \ge \lambda^+$. The clause of Definition 3.2 for a term $t$ at $\lambda$ consists of conditions on
$\lambda$ alone and of membership of $\lambda$ in further terms, the terms it \emph{requires}. The test below unfolds the
clause along every finite chain of required terms. For a term code $t$ define, with all quantifiers bounded in $V$, $\Tb$ and $\lambda$ (including
$\Lv_{\lambda^+}$ for the witnesses of Lemma 4.2):
\begin{itemize}
\item $\mathrm{Loc}(t)$: for $t = \Ad$, $\Adm(V,\lambda)$. For $t = X[b]$: $\Adm(V, \lambda)$ and $t_X \in \Tb(\eta)$ for
  cofinally many $\eta < \lambda$. For $t = H_{X,c}(Q,e)$: the same, together with $p_Q, e \in \Lv_\lambda$ and
  $\WF_V(Q,\lambda,e)$.
\item $\mathrm{Req}(t)$: $\Ad$ requires nothing. $X[b]$ requires the codes of $H_{X,c}(Q,e)$ with $c < b$, $\rho_Q$ the
  code of a $\Sigma_1$ formula, $p_Q, e \in \Lv_\lambda$, and $\WF_V(Q, \lambda, e)$ (the clause asks only for
  well-founded $e$). $H_{X,c}(Q,e)$ requires the codes of $H_{X,c}(Q,e')[c]$ with $e' \in \Lv_\lambda$ and $(\rho_Q,
  \langle e', e, p_Q\rangle) \in \mathrm{Tr}_\lambda$.
\item $\mathrm{Path}(t)$: $t \in \Lv_\lambda$ (otherwise the next condition would hold vacuously), and for every finite
  sequence $(t_0, \dots, t_n) \in \Lv_\lambda$ with $t_0 = t$ and $t_{i+1} \in \mathrm{Req}(t_i)$ for $i < n$,
  $\mathrm{Loc}(t_n)$ holds ($n = 0$ included).
\end{itemize}

\begin{statement}\head{Lemma 4.3} Let $\lambda$, $\Tb$ and $V$ be as above. For every term code $t$, $\mathrm{Path}(t)$
holds if and only if $\lambda \in t$.
\end{statement}

\begin{proof}
If $t \notin L_\lambda$, both sides fail, by Lemma 3.3(b). Let $t \in L_\lambda$. Since $\lambda$ is a limit, a finite
sequence starting with $t$ and continuing with a required $t'$ is in $L_\lambda$ exactly when its tail starting with
$t'$ is, and the required codes are in $L_\lambda$; so
\[ \mathrm{Path}(t) \iff \mathrm{Loc}(t) \wedge \forall t' \in \mathrm{Req}(t)\ \mathrm{Path}(t'). \]
By (C1), Lemma 4.2 and the choice of $\Tb$ (a member of a term at $\eta$ contains its code in $L_\eta$, Lemma 3.3(b)),
$\mathrm{Loc}(t)$ is the conjunction of the conditions of the clause of $t$ on $\lambda$ alone, and $\mathrm{Req}(t)$
is the set of terms the clause requires. So the clause reads $\lambda \in t \iff \mathrm{Loc}(t) \wedge \forall t' \in
\mathrm{Req}(t)\, (\lambda \in t')$.

If $t = H_{X,c}(Q,e)$ with $e \notin \WF(Q,\lambda)$, $\mathrm{Loc}(t)$ fails and both sides are false. Otherwise $t$ has
the measure of Lemma 3.3 at $\lambda$, and each required term has a smaller one: $H_{X,c}(Q,e)$ required by $X[b]$ has
$c < b$ and $e$ well-founded, so its measure $(c, 1, \mathrm{rank}\, e)$ is below $(b, 0, 0)$; $H_{X,c}(Q,e')[c]$
required by $H_{X,c}(Q,e)$ has measure $(c, 0, 0) < (c, 1, \mathrm{rank}\, e)$. Induction on the measure gives the
equivalence.
\end{proof}

\begin{statement}\head{Lemma 4.4 (terms are decided one admissible level up)} There is a $\Delta_0$ formula $\Phi(t,
\nu, V)$ without parameters such that for every term code $t$, every ordinal $\nu$ and every certificate $V$ through an
ordinal $\eta \ge \nu^+$: $\Phi(t, \nu, V)$ holds if and only if $\nu \in t$.
\end{statement}

\begin{proof}
\emph{The formula.} $\Phi(t,\nu,V)$ says: $\Adm(V,\nu)$, and there are $\kappa \le \eta$ with $\kappa = \nu^+$
according to $V$ and $\Tb \in \Lv^V_\kappa$ such that $\Tb$ is a function on $\nu$; for every $\lambda < \nu$ with
$\Adm(V,\lambda)$, $\Tb(\lambda)$ is the set of term codes $t' \in \Lv_\lambda$ with $\mathrm{Path}(t')$ evaluated at
$\lambda$ with the table $\Tb{\restriction}\lambda$ and the certificate $V$; for every $\lambda < \nu$ without
$\Adm(V,\lambda)$, $\Tb(\lambda) = \emptyset$ (set explicitly, since at such $\lambda$ the path of length $0$ need not
lie in $\Lv_\lambda$); and $\mathrm{Path}(t)$ holds at $\nu$ with the table $\Tb$ and $V$. All quantifiers are bounded.

\emph{Soundness.} Suppose $\Phi(t,\nu,V)$. By (C1) $V$ is the true certificate. By induction on $\lambda < \nu$,
$\Tb(\lambda)$ is the table's true row: at non-admissible $\lambda$ both are empty, since terms are subclasses of
$\Ad$; at admissible $\lambda$, $\lambda^+ \le \nu$, so $V$ supplies the data through $\lambda^+$, and Lemma 4.3 at
$\lambda$ applies, the earlier rows being true. Lemma 4.3 at $\nu$ then gives $\nu \in t$.

\emph{The table is in $L_{\nu^+}$.} Let $\nu \in \Ad$ and let $W \in L_{\nu^+}$ be the certificate through $\nu$ (C2).
For admissible $\lambda < \nu$, $\lambda^+ \le \nu$, so $W$ supplies the data for the test at $\lambda$. Inside
$L_{\nu^+}$, $\Sigma$-recursion on $\lambda < \nu$, with each row formed by $\Delta_0$-separation from $\Lv_\lambda$
(and empty at non-admissible $\lambda$), produces a function on $\nu$; by Lemma 4.3 and induction on $\lambda$ it is the
table of terms below $\nu$.

\emph{Completeness.} If $\nu \in t$, then $\nu$ is admissible and $t \in L_\nu$ (Lemma 3.3), $\kappa = \nu^+ \le
\eta$, the table is in $L_{\nu^+} = \Lv^V_\kappa$, and Lemma 4.3 at $\nu$ gives $\mathrm{Path}(t)$. If $\nu$ is not
admissible, $\nu \notin t$ and $\Adm(V,\nu)$ fails.
\end{proof}

\begin{statement}\head{Lemma 4.5 (unboundedness)} For a term $X$ let $U_X$ be the sentence
\[ \forall y\, \exists \nu, \eta, V\, [\neg\mathrm{Ord}(y) \vee (y \in \nu \wedge \Cert(V, \eta) \wedge \nu^+ \le \eta
   \text{ according to } V \wedge \Phi(t_X, \nu, V))], \]
with the constant $t_X$; it is $\Pi_2$. (a) If $t_X \in L_\beta$, $\beta > 0$ and $L_\beta \models U_X$, then $\beta \in
\Lim X$. (b) If $\mu$ is admissible and $\mu \in \Lim X$, then $t_X \in L_\mu$ and $L_\mu \models U_X$.
\end{statement}

\begin{proof}
(a) For each ordinal $y < \beta$, $L_\beta$ contains $\nu > y$ and a certificate $V$, true by (C1), with $\Phi(t_X, \nu,
V)$, so $\nu \in X$ by Lemma 4.4. Hence $X \cap \beta$ is unbounded in $\beta$, and $\beta$ is a limit.
(b) Some member of $X$ is below $\mu$, so $t_X \in L_\mu$ by Lemma 3.3(b). Given an ordinal $y < \mu$, take $\nu <
\nu'$ in $X \cap (y, \mu)$; $\nu'$ is admissible, so $\nu^+ \le \nu' < \mu$, a certificate $V$ through $\nu^+$ is in
$L_\mu$ by (C2), and $\Phi(t_X, \nu, V)$ holds by Lemma 4.4.
\end{proof}

\section{Upper bounds}

\begin{statement}\head{Lemma 5.1} Let $X$ be a term, and let $\mu$ be an admissible ordinal in $\Lim X$ that is
$\Pi_2$-reflecting on a class $Y$. Then $\mu$ is $\Pi_2$-reflecting on $Y \cap \Lim X$.
\end{statement}

\begin{proof}
Let $\varphi(q)$ be a $\Pi_2$ sentence true in $L_\mu$. By Lemma 4.5(b), $\varphi(q) \wedge U_X$ is a $\Pi_2$
sentence, with constants $q, t_X \in L_\mu$, true in $L_\mu$. A target $\beta \in Y \cap \mu$ with $q, t_X \in L_\beta$
satisfies $\varphi(q)$ and $U_X$, so $\beta \in \Lim X$ by Lemma 4.5(a).
\end{proof}

\begin{statement}\head{Theorem 5.2} Let $X$ be a term, $\mu$ an ordinal, $c, d < \omega$, and $(Q, e)$ an eligible pair
at $\mu$.
\begin{itemize}
\item[(A)] If $\mu \in \Lim X$, and $\mu$ is admissible if $c = 0$ and recursively $(c-1)$-Mahlo if $c \ge 1$, then
  $\mu \in X[c]$.
\item[(B)] If $\mu$ is $\Pi_2$-reflecting on the class of ordinals in $\Lim X$ that are admissible if $d = 0$ and
  recursively $(d-1)$-Mahlo if $d \ge 1$, then $\mu \in H_{X,d}(Q, e)$ and $\mu \in \Lim
  H_{X,d}(Q, e)$.
\end{itemize}
\end{statement}

\begin{proof}
\emph{The induction.} Give (A) at $(\mu, c)$ the measure $(\mu, 2c, 0)$ and (B) at $(\mu, d, e)$ the measure $(\mu, 2d +
1, \mathrm{rank}_\mu(e))$, ordered lexicographically, and prove both for all terms $X$ by induction on the measure. Each
appeal below is to a smaller measure: (A) at $(\mu, c)$ appeals to (B) at $(\mu, d, e)$ with $d < c$, and $2d + 1 < 2c$;
(B) at $(\mu, d, e)$ appeals to (B) at $(\mu, d, e')$ for predecessors $e'$ of $e$ (third coordinate), to (A) at $(\mu,
d)$ (second coordinate) and to (A) at levels $\beta < \mu$ (first coordinate). The measure does not involve the base
term, so the appeals may change it.

\emph{(A).} $\mu \in \Pi_2^{c+1} \subseteq \Ad$. Let $d < c$. Then $\mu$ is $\Pi_2$-reflecting on $\Pi_2^{c} \subseteq
\Pi_2^{d+1}$ (Lemma 2.1(a)), and by Lemma 5.1 on $\Pi_2^{d+1} \cap \Lim X$. By (B) at $(\mu, d, e)$, $\mu \in H_{X,d}(Q, e)$
for every eligible pair $(Q,e)$. With $\mu \in \Ad \cap \Lim X$ this is the clause of $X[c]$.

\emph{(B).} $\mu$ is $\Pi_2$-reflecting on $\Pi_2^{d+1}$, so $\mu \in \Pi_2^{d+2} \subseteq \Pi_2^{d+1}$, and $\mu$ is
admissible. Reflecting ``$y$ exists'' for each $y < \mu$ gives members of $\Lim X$ cofinal in $\mu$, so $\mu \in \Lim
X$. Write $H(e') = H_{X,d}(Q, e')$.

\emph{Membership.} Let $e'$ be a predecessor of $e$ at $\mu$. By (B) at $(\mu, d, e')$, $\mu \in \Lim H(e')$, so by (A)
at $(\mu, d)$ with base term $H(e')$, $\mu \in H(e')[d]$. With $\mu \in \Ad \cap \Lim X$ and $p_Q, e \in L_\mu$, $\mu \in
H(e)$.

\emph{Limit.} Write $\rho_Q = \exists i\, \rho_{Q,0}$ with $\rho_{Q,0} \in \Delta_0$, and let $\Theta$ be the sentence
\[ \begin{aligned} & \forall e', z, i\ \exists \nu, \eta, V, t\, [\, \neg\mathrm{Ord}(z) \vee \neg\rho_{Q,0}(e', e,
   p_Q, i) \vee (z \in \nu \wedge \Cert(V,\eta) \\ & \qquad \wedge \nu^+ \le \eta \text{ according to } V \wedge t =
   t_{H(e')} \wedge \Phi(t, \nu, V))\,], \end{aligned} \]
where $t_{H(e')}$ is the code of $H(e')$, formed from $e'$ and the constants by tuple formation. It is $\Pi_2$, with
constants $e, p_Q, t_X \in L_\mu$. It holds in $L_\mu$: for a predecessor $e'$ of $e$ and $z < \mu$, take $\nu < \nu'$
in $H(e') \cap (z, \mu)$, which exist because $\mu \in \Lim H(e')$; $\nu'$ is admissible, so $\nu^+ < \mu$, (C2) gives a
certificate through $\nu^+$ in $L_\mu$, and Lemma 4.4 gives $\Phi$. Given an ordinal $y < \mu$, reflect $\Theta$, with
$y$ added to its constants, onto $\Pi_2^{d+1} \cap \Lim X$: some $\beta$ with $y < \beta < \mu$ is in $\Pi_2^{d+1} \cap \Lim
X$ and $L_\beta \models \Theta$. Then $\beta$ is admissible, and $e \in \WF(Q, \beta)$ by (C5). For a predecessor $e'$
of $e$ at $\beta$, $\Theta$ in $L_\beta$, (C1) and Lemma 4.4 give members of $H(e')$ cofinal in $\beta$, so $\beta \in
\Lim H(e')$, and (A) at $(\beta, d)$ gives $\beta \in H(e')[d]$. Hence $\beta \in H(e)$, and $\mu \in \Lim H(e)$.
\end{proof}

\begin{statement}\head{Lemma 5.3} Let $X$ be a term and $\gamma$ an admissible ordinal.
\begin{itemize}
\item[(a)] If $c < \omega$ and $\gamma$ is a limit of recursively $c$-Mahlo members of $\Lim X$, then $\gamma \in H_{X,c}(Q, e)$ for every
  eligible pair $(Q,e)$ at $\gamma$.
\item[(b)] If $k < \omega$ and $\gamma$ is a limit of members of $\Lim X$ that are admissible if $k = 0$ and recursively $(k-1)$-Mahlo if $k \ge 1$, then
  $\gamma \in X[k]$.
\end{itemize}
\end{statement}

\begin{proof}
In both parts $\gamma \in \Lim X$, since a limit of members of $\Lim X$ is in $\Lim X$. We prove (b) by induction on
$k$, together with (a) for every $c < k$.

(b) for $k = 0$: $X[0] = \Ad \cap \Lim X$.

(a) for $c$, assuming (b) for $c$ and every term. $\gamma \in \Ad \cap \Lim X$ and $p_Q, e \in L_\gamma$. Let $e'$ be a
predecessor of $e$ at $\gamma$ and write $H(e') = H_{X,c}(Q,e')$. Points $\mu \in \Pi_2^{c+2} \cap \Lim X$ below $\gamma$
with $p_Q, e' \in L_\mu$ are cofinal in $\gamma$, and for each of them $(Q, e')$ is an eligible pair at $\mu$ by (C5).
Such a $\mu$ is admissible and $\Pi_2$-reflecting on $\Pi_2^{c+1}$, hence, by Lemma 5.1, on $\Pi_2^{c+1} \cap \Lim X$;
Theorem 5.2(B) gives $\mu \in H(e') \cap \Lim H(e')$, and $\mu \in \Pi_2^{c+1}$ by Lemma 2.1(a). So $\gamma \in
\Lim(\Pi_2^{c+1} \cap \Lim H(e'))$, and (b) for $c$ with base term $H(e')$ gives $\gamma \in H(e')[c]$. Hence $\gamma \in
H_{X,c}(Q,e)$.

(b) for $k+1$, assuming (b) for every $k' \le k$. $\gamma \in \Ad \cap \Lim X$. Let $c \le k$. Since $\Pi_2^{k+2} \subseteq
\Pi_2^{c+2}$, $\gamma \in \Lim(\Pi_2^{c+2} \cap \Lim X)$, and (a) for $c$ gives $\gamma \in H_{X,c}(Q,e)$ for every eligible
pair. So $\gamma \in X[k+1]$.
\end{proof}

\section{The lower bound}

The search tree is that of [A, Definition 5.3] with two changes: a positive answer must lie in $\Pi_2^{k+1}$ rather than in
$\Ad$, and each record carries a point of $X$, so that a branch has its supremum in $\Lim X$.

\begin{statement}\head{Definition 6.1 (the tree $S$)} Fix $k < \omega$, a term $X$ and an ordinal $\ell$.

\emph{Tasks.} Task $0$ is empty. For $m \ge 1$ write $m = 2^i(2n+1)$; task $m$ is $\mathrm{Refl}(i, n)$. Each pair $(i,
n)$ is task $m$ for exactly one $m$, and then $i < m$.

\emph{Nodes.} A node of depth $m \ge 0$ is a sequence of records $r_0, \dots, r_m$. Record $r_0$ is an ordinal $a_0 >
\max(\ell, \omega)$. For $1 \le j \le m$, record $r_j$ is a tuple $(a_j, f_{j-1}, z_j, x_j)$ with:
\begin{itemize}
\item an ordinal $a_j > a_{j-1}$, and a surjection $f_{j-1} \in L_{a_j}$ from $\omega$ onto $L_{a_{j-1}}$;
\item a point $z_j \in X$ with $a_{j-1} < z_j$ and $z_j^+ < a_j$;
\item the answer $x_j \in L_{a_j}$ to task $j-1$.
\end{itemize}
The ordinals $a_j$ are the \emph{stages} of the node. Task $j - 1 = \mathrm{Refl}(i,n)$ has $i < j - 1$, so $f_i$ is in
the node. If $f_i(n)$ is a pair $(\psi, q)$ with $\psi$ the code of a $\Pi_2$ formula $\forall u\, \exists v\,
\theta(q,u,v)$, the task is \emph{about} $(\psi, q)$, and its answer is $(+, b)$ with $b$ admissible if $k = 0$ and recursively $(k-1)$-Mahlo if $k \ge 1$, $b < a_j$, $q \in
L_b$ and $L_b \models \psi(q)$, or $(-, u)$, a \emph{promise} that $\forall v\, \neg\theta(q, u, v)$. Otherwise the
task is empty and its answer is $\emptyset$. \emph{Node condition:} every promise $(-, u)$ of a node of depth $m$
satisfies $\neg\theta(q, u, v)$ for all $v \in L_{a_m}$.

\emph{Edges.} $S$ has the parameter $(\ell, k, t_X)$. The predecessors of a node of depth $m$ are its \emph{children},
the nodes of depth $m+1$ extending it, and the root $r = \emptyset$ has the nodes of depth $0$ as predecessors. The
\emph{next task} of a node of depth $m$ is task $m$, answered in record $r_{m+1}$ of a child.
\end{statement}

For example, a node of depth 1 is $(a_0, (a_1, f_0, z_1, \emptyset))$, since task $0$ is empty; its next task is task
$1 = \mathrm{Refl}(0,0)$, about $f_0(0)$ if that is a pair of the right form. Along an infinite branch the $f_j$
enumerate $L_\lambda$ for $\lambda = \sup_j a_j$, every $\Pi_2$ sentence with constant in $L_\lambda$ is eventually
answered, and the points $z_j$ make $\lambda$ a limit of $X$ (Lemma 6.3).

\begin{statement}\head{Lemma 6.2} $S$ is a $\Sigma_1$ presentation. At every limit $\mu > \omega$ containing its
parameter, its edges are edges of the tree of Definition 6.1, and at admissible $\mu$ containing the parameter every edge $(s', s)$
of that tree with $s' \in L_\mu$ is an edge at $\mu$. For $\max(\ell,\omega) < a_0 < \mu$, the node $(a_0)$ is a predecessor of $r$
at every limit $\mu > a_0$ containing the parameter.
\end{statement}

\begin{proof}
The edge formula is the disjunction of a root-edge clause and a positive-depth clause. The root-edge clause, ``$s =
r$ and $s' = (a_0)$ with $a_0$ an ordinal above $\max(\ell, \omega)$'', is bounded. The positive-depth clause, ``$s'$
is a node extending $s$ by one record'', is $\exists V\, [\Cert(V, a) \wedge B]$, where $a$ is the last stage of $s'$ and
$B$ is bounded: $b \in \Pi_2^{k+1}$ is $\chi(V{\restriction}(b+1), b, k+1)$ (Lemma 4.1); $z_j^+ < a_j$ is $\exists
\kappa < a_j\, (z_j < \kappa \wedge \Adm(V, \kappa))$; given this, $V{\restriction}(a_j+1)$ is a certificate through an
ordinal $\ge z_j^+$, so $z_j \in X$ is $\Phi(t_X, z_j, V{\restriction}(a_j+1))$ (Lemma 4.4); $L_b \models \psi(q)$
and the node condition are read from $\mathrm{Tr}_b$ and $\mathrm{Tr}_{a_m}$, and the tasks from $\mathrm{Tr}_\omega$.
So the edge relation is $\Sigma_1$, with prefix $\exists V$. By (C1) any such $V$ is true, so edges at $\mu$ are edges
of the tree; at admissible $\mu$, (C2) supplies $V$ for an edge with $s' \in L_\mu$. The root-edge clause holds at every limit $\mu > a_0$ containing the parameter, with no certificate, which gives
the last assertion.
\end{proof}

\begin{statement}\head{Lemma 6.3 (branches)} If the root is not in $\WF(S, \mu)$ at a limit $\mu > \omega$ containing the
parameter, then some recursively $k$-Mahlo $\lambda \in (\ell, \mu]$ lies in $\Lim X$.
\end{statement}

\begin{proof}
Take an infinite sequence $r, s_0, s_1, \dots$ with each term a predecessor of the previous one at $\mu$; by Lemma 6.2,
$s_m$ is a node of depth $m$ and the $s_m$ extend one another. Their stages $a_0 < a_1 < \cdots$ are below $\mu$; let
$\lambda = \sup_j a_j$, a limit above $\omega$ with $\ell < \lambda \le \mu$. Since $a_{j-1} < z_j < a_j$ with $z_j \in
X$, $\lambda \in \Lim X$, and $L_\lambda$ is the union of the ranges of the $f_j$. Let $\psi(q) = \forall u\, \exists
v\, \theta$ with $q \in L_\lambda$ be true in $L_\lambda$. Then $(\psi, q) = f_i(n)$ for some $i, n$, and
$\mathrm{Refl}(i,n)$ is task $m$ for some $m > i$, answered by $x_{m+1} \in L_{a_{m+1}}$. A promise $(-, u)$ would have
$u \in L_{a_{m+1}}$, and the node conditions of the $s_{m'}$, $m' \ge m+1$, would give $\neg\theta(q, u, v)$ for every
$v \in L_\lambda$, so $\psi(q)$ would be false in $L_\lambda$. So the answer is $(+, b)$ with $b \in \Pi_2^{k+1} \cap
\lambda$, $q \in L_b$ and $L_b \models \psi(q)$. By block contraction at $\lambda$ and at $b$, $\lambda$ is
$\Pi_2$-reflecting on $\Pi_2^{k+1}$, that is, $\lambda \in \Pi_2^{k+2}$.
\end{proof}

For $\beta \in \Pi_2^{k+1} \cap \Lim X$ with $\ell < \beta$, $\beta$ at most the least (+1)-stable ordinal, and the
parameter in $L_\beta$, and a node $s \in
L_\beta$: $s$ is \emph{correct at $\beta$} if each of its promises holds in $L_\beta$ (these are $\Pi_1$ statements). A
task about $(\psi, q)$ is \emph{bad at $\beta$} if $L_\beta \models \psi(q)$ and $L_b \nvDash \psi(q)$ for every $b \in
\Pi_2^{k+1} \cap \beta$ with $q \in L_b$.

\begin{statement}\head{Lemma 6.4 (extension)} Let $\beta$ be as above and $s \in L_\beta$ a node of depth $m$, correct
at $\beta$, whose next task is not bad at $\beta$. Then $s$ has a child $s' \in L_\beta$ that is a predecessor of $s$ at
$\beta$ and correct at $\beta$, and there is $w \in L_\beta$ with this property: for every $\beta' < \beta$ in
$\Lim X$ that is admissible if $k = 0$ and recursively $(k-1)$-Mahlo if $k \ge 1$, with $s', w \in L_{\beta'}$, the node $s'$ is correct at $\beta'$ and its next task is not bad
at $\beta'$.
\end{statement}

The set $w$ records how the next task of $s'$ is settled at $\beta$, so that it stays settled at smaller levels that
contain $w$.

\begin{proof}
\emph{The answer.} Let task $m$ be about $(\psi, q)$, $q \in L_{a_m}$. Answer $\emptyset$ if it is empty; $(-, u)$ with
$u \in L_\beta$ and $L_\beta \models \forall v\, \neg\theta(q, u, v)$ if $\psi(q)$ is false in $L_\beta$; and $(+, b)$
with $b \in \Pi_2^{k+1} \cap \beta$, $q \in L_b$ and $L_b \models \psi(q)$ if it is true, which exists because the task is
not bad.

\emph{The new record.} $L_{a_m}$ is countable in $L_\beta$, since $\beta$ is admissible and at most the least (+1)-stable ordinal, which
gives
$f_m \in L_\beta$. Since $\beta \in \Lim X$, there are $z_{m+1} < z'$ in $X \cap (a_m, \beta)$; $z'$ is admissible, so
$z_{m+1}^+ \le z' < \beta$. Choose $a_{m+1} < \beta$ above $z_{m+1}^+$ with $f_m$ and the answer in $L_{a_{m+1}}$. The
promises of $s$ and the new one hold in $L_\beta$, hence in $L_{a_{m+1}}$, since $\Pi_1$ sentences pass to transitive
subsets containing their constants. So $s'$ is a child of $s$; it is a predecessor of $s$ at $\beta$ by Lemma 6.2, and
correct at $\beta$.

\emph{The set $w$.} Task $m+1$ is empty or about some $(\psi', q')$ with $q' \in L_{a_m}$; write $\psi' = \forall u\,
\exists v\, \theta'$. Let $w$ be $u$ with $L_\beta \models \forall v\, \neg\theta'(q', u, v)$ if $\psi'(q')$ is false in
$L_\beta$; $b \in \Pi_2^{k+1} \cap \beta$ with $q' \in L_b$ and $L_b \models \psi'(q')$ if it is true and not bad; and
$\emptyset$ otherwise.

\emph{Passing down.} Let $\beta' < \beta$ be in $\Pi_2^{k+1} \cap \Lim X$ with $s', w \in L_{\beta'}$. Then $\ell < a_0 <
\beta' < \beta$, so $\beta'$ is at most the least (+1)-stable ordinal, and the parameter is in $L_{\beta'}$ ($t_X$ by Lemma 3.3(b), since $\beta' \in \Lim X$), so
correctness and badness are defined at $\beta'$. The promises of $s'$ are $\Pi_1$ with constants in $L_{\beta'}$, so
$s'$ is correct at $\beta'$. Task $m+1$ is not bad at $\beta'$: if it is empty there is nothing to show; if $\psi'(q')$
was false at $\beta$ with witness $u \in L_{\beta'}$, it is false at $\beta'$; if it was true and not bad, the witness $b
< \beta'$ shows it is not bad at $\beta'$; if it was bad at $\beta$, then $\beta'$ is an element of $\Pi_2^{k+1} \cap
\beta$ with $q' \in L_{\beta'}$, so $\psi'(q')$ is false at $\beta'$.
\end{proof}

\begin{statement}\head{Theorem 6.5 (lower bound)} Let $k < \omega$, $X$ a term and $R$ the class of members of
$\Lim X$ that are admissible if $k = 0$ and recursively $(k-1)$-Mahlo if $k \ge 1$. For
every admissible $\gamma$ at most the least (+1)-stable ordinal, if $\gamma \in X[k]$, then $\gamma \in R \cup \Lim R$.
\end{statement}

\begin{proof}
Induction on $k$, for all terms $X$ at once.

\emph{$k = 0$.} $X[0] = \Ad \cap \Lim X = \Pi_2 \cap \Lim X$, so $\gamma \in R$.

\emph{From $k$ to $k+1$.} Let $\gamma$ be admissible, at most the least (+1)-stable ordinal, with $\gamma \in
X[k+1]$, let $\ell < \gamma$, and
suppose that no point of $R' = \Pi_2^{k+2} \cap \Lim X$ lies in $(\ell, \gamma]$. Let $S = S^{k,X}_\ell$ be the tree of
Definition 6.1; its parameter is in $L_\gamma$, since $\gamma \in \Lim X$ (Lemma 3.3(b)). For a node $s$ write $H(s) =
H_{X,k}(S, s)$; each $H(s)$ is a term, and $H(s) \subseteq \Lim X$ by its clause.

\emph{The root.} By Lemma 6.3 and the assumption, the root $r$ is in $\WF(S, \gamma)$. As $k < k+1$ and $(S, r)$ is an
eligible pair at $\gamma$, $\gamma \in H(r)$. The assumption on $(\ell, \gamma]$ is used only here. The rest of the
proof descends from $\gamma$ through smaller members of the classes $H(s)$, using the induction hypothesis to keep each
level in $\Pi_2^{k+1}$.

\emph{Claim.} Let $\beta \in \Pi_2^{k+1} \cap \Lim X$ with $\ell < \beta$, $\beta$ at most the least (+1)-stable
ordinal, and the parameter in $L_\beta$, and let
$s \in L_\beta$ be a node correct at $\beta$ whose next task is not bad at $\beta$. Then $\beta \notin \Lim H(s)$.

Suppose not, and let $\beta$ be least, with such an $s$.
\begin{itemize}
\item Lemma 6.4 gives a child $s' \in L_\beta$ of $s$ and a set $w \in L_\beta$.
\item Since $\beta \in \Lim H(s)$, some $\beta_0 < \beta$ in $H(s)$ has $s', w \in L_{\beta_0}$. Then $\beta_0$ is
  admissible, so $s'$ is a predecessor of $s$ at $\beta_0$ by Lemma 6.2. The clause of $H(s)$ at $\beta_0$ gives $s \in
  \WF(S, \beta_0)$, hence $s' \in \WF(S, \beta_0)$, and $\beta_0 \in H(s')[k]$.
\item The induction hypothesis, applied to the term $H(s')$ at the admissible $\beta_0 < \beta$, gives $\beta_0 \in
  R_{s'} \cup \Lim R_{s'}$ with $R_{s'} = \Pi_2^{k+1} \cap \Lim H(s')$. Let $\beta_1 = \beta_0$ if $\beta_0 \in R_{s'}$, and
  otherwise let $\beta_1 < \beta_0$ be in $R_{s'}$ with $s', w \in L_{\beta_1}$.
\item Then $\beta_1 \in \Pi_2^{k+1}$, and $\beta_1 \in \Lim H(s') \subseteq \Lim X$. By Lemma 6.4, applied with $\beta' =
  \beta_1$, $s'$ is correct at $\beta_1$ and its next task is not bad at $\beta_1$; and $\ell < \beta_1 < \beta$, with
  the parameter in $L_{\beta_1}$.
\item So $(\beta_1, s')$ satisfies the hypotheses of the claim, and $\beta_1 \in \Lim H(s')$, contradicting the
  minimality of $\beta$.
\end{itemize}

\emph{Conclusion.} Choose $a_0$ with $\max(\ell, \omega) < a_0 < \gamma$ and let $s_0 = (a_0)$, a predecessor of $r$ at
$\gamma$ (Lemma 6.2). The clause of $H(r)$ at $\gamma$ gives $\gamma \in H(s_0)[k]$, and the induction hypothesis gives
some $\beta_1 \in \Pi_2^{k+1} \cap \Lim H(s_0)$ with $a_0 < \beta_1 \le \gamma$. Then $\beta_1 \in \Lim X$, $\ell < \beta_1
\le \gamma$, the parameter and $s_0$ are in $L_{\beta_1}$, $s_0$ has no promises, and its next task, task $0$, is empty.
The claim gives $\beta_1 \notin \Lim H(s_0)$, a contradiction. So every $\ell < \gamma$ has a point of $R'$ in $(\ell,
\gamma]$, that is, $\gamma \in R' \cup \Lim R'$.
\end{proof}

\section{Results}

\begin{statement}\head{Theorem 7.1} Let $k < \omega$, $X$ a term and $R$ the class of members of $\Lim X$ that are
admissible if $k = 0$ and recursively $(k-1)$-Mahlo if $k \ge 1$. An admissible $\gamma$
at most the least (+1)-stable ordinal is in $X[k]$ if and only if $\gamma \in R \cup \Lim R$.
\end{statement}

\begin{proof}
The forward direction is Theorem 6.5. Conversely, $R \subseteq X[k]$ by Theorem 5.2(A), and an admissible $\gamma \in
\Lim R$ is in $X[k]$ by Lemma 5.3(b).
\end{proof}

\begin{statement}\head{Theorem 7.2 (main theorem)} Let $k \ge 0$. An admissible ordinal $\gamma$ that is at most
the least (+1)-stable ordinal is recursively inaccessible of degree $\delta_\gamma^{k+1}$ if and only if $\gamma$ is
recursively $k$-Mahlo or a limit of recursively $k$-Mahlo ordinals. The least ordinal $\alpha$ that is recursively
inaccessible of degree $\delta_\alpha^{k+1}$ is the least recursively $k$-Mahlo ordinal.
\end{statement}

\begin{proof}
Take $X = \Ad$ and $k + 1$ in place of $k$ in Theorem 7.1. Since $\Pi_2^{k+2} \subseteq \Lim \Ad$ (Lemma 2.1), $R =
\Pi_2^{k+2}$. By Lemma 2.1(b) the least element of $\Pi_2^{k+2}$ exists and is below $\pi_3$, which is at most the least
(+1)-stable ordinal, so it is in $\Ad[k+1]$. Every member of $\Ad[k+1]$ is admissible, and one that is at most the
least (+1)-stable ordinal is in $\Pi_2^{k+2} \cup \Lim
\Pi_2^{k+2}$, so it is at least the least element of $\Pi_2^{k+2}$.
\end{proof}

\begin{statement}\head{Corollary 7.3 (degrees $\delta_\alpha + 1$ and $\delta_\alpha \cdot 2$)}
\begin{itemize}
\item[(a)] The least ordinal recursively inaccessible of degree $\delta_\alpha + 1$ is the least admissible limit of
  recursively Mahlo ordinals.
\item[(b)] The least ordinal recursively inaccessible of degree $\delta_\alpha \cdot 2$ is the least recursively Mahlo
  limit of recursively Mahlo ordinals, and it is below the least recursively 1-Mahlo ordinal.
\end{itemize}
\end{statement}

\begin{proof}
For $\mu$ at most the least (+1)-stable ordinal, $D \cap \mu = (\Pi_2^2 \cup (\Ad \cap \Lim \Pi_2^2)) \cap \mu$ [A, Corollary 5.9], so $\mu \in \Lim D$
if and only if $\mu \in \Lim \Pi_2^2$. By [A, Corollary 5.9 and its proof], the least admissible limit of $\Pi_2^2$ and a
recursively Mahlo limit of recursively Mahlo ordinals exist below the least recursively 1-Mahlo ordinal, which is
below $\pi_3$, hence below the least (+1)-stable ordinal (Lemma 2.1).

(a) follows, since the ordinals of degree $\delta_\alpha + 1$ are those in $\Ad \cap \Lim D$.

(b) The class of degree $\delta_\alpha \cdot 2$ is $D[1]$ (Lemma 3.5). By Theorem 7.1 with $k = 1$ and $X = D$, an
admissible $\gamma$ at most the least (+1)-stable ordinal is in it if and only if $\gamma \in R \cup \Lim R$ with $R =
\Pi_2^2 \cap \Lim D$, and by the first paragraph $R$ and $\Pi_2^2 \cap \Lim \Pi_2^2$ have the same members up to and
including the least (+1)-stable ordinal. Every element of $\Lim R$ is above an
element of $R$, so the least element of $D[1]$ is the least element of $\Pi_2^2 \cap \Lim \Pi_2^2$.
\end{proof}

\begin{statement}\head{Proposition 7.4 (degree $\delta_\alpha^\omega$)} An admissible $\gamma$ at most the least (+1)-stable
ordinal is recursively inaccessible of degree $\delta_\gamma^\omega$ if and only if $\gamma$ is a limit of
recursively $k$-Mahlo ordinals for every $k < \omega$. The least ordinal of degree $\delta_\alpha^\omega$ is the least
admissible ordinal that is a limit of recursively $k$-Mahlo ordinals for every $k < \omega$, and it is below $m_\omega$,
the least recursively $\omega$-Mahlo ordinal.
\end{statement}

\begin{proof}
\emph{Necessity.} If $\gamma \in \Ad[\omega]$ then $\gamma \in \Ad[k+1]$ for every $k < \omega$: the clause of
$\Ad[\omega]$ asks for the classes $H_{\Ad,c}(Q,e)$ with $c < \omega$, and that of $\Ad[k+1]$ for those with $c \le k$.
By Theorem 7.2, $\gamma$ is recursively $k$-Mahlo or a limit of such ordinals, hence, by Lemma 2.1(a), a limit of
$\Pi_2^{k+1}$. So $\gamma$ is a limit of every $\Pi_2^{k+1}$, that is, of recursively $k$-Mahlo ordinals for every
$k$.

\emph{Sufficiency} (for every admissible $\gamma$). $\gamma \in \Ad \cap \Lim \Ad$. For $c < \omega$, $\gamma \in \Lim
\Pi_2^{c+2}$ ${}= \Lim(\Pi_2^{c+2} \cap \Lim \Ad)$, so Lemma 5.3(a) gives $\gamma \in H_{\Ad,c}(Q,e)$ for every eligible pair
$(Q,e)$. So $\gamma \in \Ad[\omega]$.

\emph{The least element.} Write $\tau_\omega$ for the least admissible ordinal that is a limit of recursively
$k$-Mahlo ordinals for every $k$. Then $\tau_\omega \le m_\omega$, which is below the least (+1)-stable ordinal (Lemma 2.1), so $\tau_\omega$ has the degree, and every
smaller admissible ordinal fails the characterisation.

\emph{$\tau_\omega < m_\omega$.} Let $\Psi$ be the sentence
\[ G \wedge \forall k\, \forall y\, \exists \nu\, \exists W\, [\neg(k \in \omega) \vee \neg\mathrm{Ord}(y) \vee (y \in
   \nu \wedge \Cert(W, \nu) \wedge \chi(W, \nu, k))], \]
with $\chi$ from Lemma 4.1 and the constant $\omega$. Merging the blocks of $G$, it is $\Pi_2$. It holds in
$L_{m_\omega}$: $m_\omega$ is a limit of every $\Pi_2^{k+1}$ (Lemma 2.1), and (C2) supplies the certificates. As $m_\omega \in
\Pi_2^2$, it reflects $\Psi$ onto $\Ad$: some admissible $\beta < m_\omega$ has $L_\beta \models \Psi$. For each $k <
\omega$ and $y < \beta$, $L_\beta$ contains $\nu > y$ and a certificate through $\nu$, true by (C1), with $\chi(W, \nu,
k)$, so by Lemma 4.1 $\nu \in \Pi_2^{k}$ when $k \ge 1$. So $\beta$ is an admissible limit of every
$\Pi_2^{k+1}$, and $\tau_\omega \le \beta <
m_\omega$.
\end{proof}

\section{Where the statements fail, and open questions}

\begin{itemize}
\item \emph{Exponent $0$.} The pattern of Theorem 7.2 does not extend to exponent $0$: $\Ad[0] = \Ad \cap \Lim \Ad$,
  whose least element is the least recursively inaccessible ordinal, while the matching class, that of the admissible ordinals, has least
  element $\omega_1^{CK}$, which is not a limit of admissible ordinals. Theorem 7.1 holds at $k = 0$, with $R = \Ad \cap
  \Lim X$.
\item \emph{Degree $\delta_\alpha \cdot 2$.} It does not give recursive 1-Mahloness (Corollary 7.3(b)). By Theorem 7.1
  with $k = 1$ and Lemma 3.5(b), $\delta_\alpha$ steps of ``admissible limit of'' over a base term $X$ give,
  at admissible levels up to and including the least (+1)-stable ordinal, the class $R \cup \Lim R$ with $R$ the recursively Mahlo members of $\Lim X$, and so use $X$ only through $\Lim X$; exponent $2$ iterates $Y \mapsto Y[1]$, and its least point is
  the least recursively 1-Mahlo ordinal by Theorem 7.2.
\item \emph{Exponent $\omega$.} ``The least ordinal of degree $\delta_\alpha^{k+1}$ is the least recursively
  $k$-Mahlo ordinal'' holds for $k < \omega$, and its continuation fails at $\omega$: the least ordinal of degree
  $\delta_\alpha^\omega$ is below $m_\omega$, the least recursively $\omega$-Mahlo ordinal (Proposition 7.4). At
  exponent $\omega$ the conditions ``recursively $k$-Mahlo or a limit of recursively $k$-Mahlo ordinals'' are
  intersected over $k$, while recursive $\omega$-Mahloness intersects the conditions ``recursively $k$-Mahlo''.
\item \emph{Parameters.} Theorem 5.2 reflects sentences that mention the tests of Lemmas 4.1 and 4.4, and Lemma 6.2
  uses the fixed tests in the edge formula of a $\Sigma_1$ presentation with parameter $(\ell, k, t_X)$, so the constants of these tests must lie in the levels
  concerned; this holds because the tests have no parameters and the constants of a term lie below its members (Lemma
  3.3(b)). If the class $\Ad$ in Definition 3.1 is replaced by a class tested only with a constant outside the level,
  the upper bound fails: for the least recursively Mahlo ordinal $\mu$, take the class consisting of $\mu$ and an
  $\omega$-sequence of admissible ordinals cofinal in $\mu$, chosen outside $L_\mu$ so that every true $\Pi_2$ sentence
  of $L_\mu$ holds at one of them; then $\mu$ is $\Pi_2$-reflecting on this class, but no member below $\mu$ is a limit
  of the class, so two steps of ``in the class and a limit of'' already exclude $\mu$.
\end{itemize}

\textbf{Open questions.}
\begin{enumerate}
\item Allowing every $\Sigma_1$ presentation in place of $P_\omega$ to supply the exponent in Definition 3.2 defines,
  at limits $\alpha > \omega$, degree $\delta_\alpha^{\delta_\alpha}$. Is its least element the least $\Pi_3$-reflecting ordinal $\pi_3$, as the third
  question of [A] asks? We do not treat this degree here.
\item What are the least elements for the exponent $\omega + 1$ and for exponents given by other nodes of $\Sigma_1$
  presentations?
\item Does the forward direction of Theorem 7.2 hold above the least (+1)-stable ordinal? The proof uses that bound
  only for the
  countability in Lemma 6.4.
\end{enumerate}

\section{Verification status}

\begin{center}\begin{tabular}{p{0.24\linewidth}p{0.70\linewidth}}
\toprule
Part & Status \\
\midrule
All numbered results & Written proofs, not checked by a person. The upper bound (Theorem 5.2) adapts known arguments. \\
AI check & A referee from a different model family checked the written arguments in Sections 2--8. \\
Lean & None. \\
References & Results of [A] and [E] are used as stated there. Richter--Aczel is cited from the 1973 Oslo preprint. The
  locators in Arai's Lectures (arXiv v1, Section 4) and in Rathjen's 1993 preprint (Remark 4.9) were checked against
  those texts. \\
\bottomrule
\end{tabular}\end{center}

\textbf{Novelty.} Not established; Section 1 compares the nearest results.

The passages a specialist is most likely to question are the bounded test of Lemmas 4.3 and 4.4, in particular that the
table of terms lies in $L_{\nu^+}$ itself; the claim in the proof of Theorem 6.5, in particular the choice of
$\beta_1$ that keeps $s'$ and $w$ and lies in $\Pi_2^{k+1}$; and the extension step of Lemma 6.4, with the point
$z_{m+1}$ and the bound $z_{m+1}^+ < a_{m+1}$.

\section*{Appendix A: which result uses which}

\begin{center}\begin{tabular}{p{0.26\linewidth}p{0.68\linewidth}}
\toprule
Result & Uses \\
\midrule
Lemma 2.1 & [A, Theorem 4.3]; [E, Definition 3.1, Example 3.2(1)] \\
Lemmas 3.3, 3.5 & Definitions 3.1--3.4; [E, (C5)] \\
Lemma 4.1 & [E, (C1)--(C4)]; Richter--Aczel 1.8 \\
Lemmas 4.2--4.4 & Lemma 3.3; [E, (C1), (C2)]; $\Sigma$-recursion, $\Delta_0$-separation and $\Delta_0$-collection \\
Lemma 4.5 & Lemmas 3.3, 4.4; [E, (C1), (C2)] \\
Lemma 5.1 & Lemma 4.5 \\
Theorem 5.2 & Lemmas 2.1, 4.4, 5.1; [E, (C1), (C2), (C5)] \\
Lemma 5.3 & Theorem 5.2; Lemmas 2.1, 5.1; [E, (C5)] \\
Lemmas 6.2--6.4 & Lemmas 3.3, 4.1, 4.4; [E, (C1), (C2)]; [E, Lemma 7.2] for the countability in Lemma 6.4 \\
Theorem 6.5 & Lemmas 3.3, 6.2--6.4 \\
Theorems 7.1, 7.2 & Theorem 6.5, Theorem 5.2, Lemmas 2.1, 5.3 \\
Corollary 7.3 & Theorem 7.1, Lemmas 2.1, 3.5; [A, Corollary 5.9] \\
Proposition 7.4 & Theorem 7.2, Lemmas 2.1, 4.1, 5.3; [E, (C1), (C2)] \\
\bottomrule
\end{tabular}\end{center}

\section*{References}

\begin{itemize}
\item $[A]$ \href{https://waffle3z.github.io/math/degree-reflection.html}{The least ordinal that is $\Pi_n$-reflecting of
  degree $\delta_\alpha$ is the least $\Pi_{n+1}$-reflecting ordinal}, waffle3z.github.io/math.
\item $[E]$ \href{https://waffle3z.github.io/math/greatly-mahlo-pi3.html}{The least $\alpha^+$-recursively Mahlo ordinal
  is the least $\Pi_3$-reflecting ordinal}, waffle3z.github.io/math.
\item J. P. Aguilera, The order of reflection, arXiv:1906.11769 (Section 2).
\item T. Arai, Iterating the recursively Mahlo operations, arXiv:1005.1987 (Sections 2--3).
\item T. Arai, Lectures on Ordinal Analysis, arXiv:2304.00246v1 (Section 4).
\item J. Barwise, \emph{Admissible Sets and Structures}, Springer, 1975 (Corollary I.8.6, Theorem I.6.4).
\item M. Rathjen, Proof theory of reflection, preprint, 1993 (Remark 4.9); published in \emph{Annals of Pure and Applied
  Logic} 68 (1994), 181--224.
\item W. Richter and P. Aczel, Inductive definitions and reflecting properties of admissible ordinals, in
  \emph{Generalized Recursion Theory}, North-Holland, 1974, 301--381 (Oslo preprint, 1973: Theorems 1.8--1.11).
\end{itemize}

\end{document}
